
\documentclass{article}
\usepackage{colortbl}
\usepackage{makecell}
\usepackage{multirow}
\usepackage{supertabular}

\begin{document}

\newcounter{utterance}

\centering \large Interaction Transcript for game `eqbench', experiment `eq\_bench', episode 68 with reflexion{-}llama{-}3.1{-}8b{-}instruct.
\vspace{24pt}

{ \footnotesize  \setcounter{utterance}{1}
\setlength{\tabcolsep}{0pt}
\begin{supertabular}{c@{$\;$}|p{.15\linewidth}@{}p{.15\linewidth}p{.15\linewidth}p{.15\linewidth}p{.15\linewidth}p{.15\linewidth}}
    \# & \multicolumn{2}{c}{Player} && \multicolumn{2}{c}{Game Master} \\
    \hline

    \theutterance \stepcounter{utterance}  
    & & & \multicolumn{4}{p{0.6\linewidth}}{
        \cellcolor[rgb]{0.9,0.9,0.9}{
            \makecell[{{p{\linewidth}}}]{
                \texttt{\tiny{[P1$\langle$GM]}}
                \texttt{You have attempted this same task before and did not succeed. These are the plans you wrote after those attempts.} \\
\texttt{Trial \#1: New plan: } \\
\\ 
\texttt{Upon re{-}examining the dialogue, I realized that I focused too much on the surface{-}level emotions expressed by Felix and Vanessa, but didn't consider the underlying dynamics of their relationship and the severity of Felix's actions. } \\
\\ 
\texttt{To improve my prediction, I will pay closer attention to the following:} \\
\\ 
\texttt{1. The extent of Felix's betrayal: I will consider the impact of Felix's actions on Vanessa's reputation, career, and trust in their friendship.} \\
\texttt{2. Vanessa's emotional state: I will analyze her reactions and language to gauge the depth of her hurt, anger, and disappointment.} \\
\texttt{3. Felix's lack of accountability: I will assess his responses to see if he truly understands the gravity of his mistakes and takes responsibility for his actions.} \\
\\ 
\texttt{With these factors in mind, I will reassess the emotions Felix is likely to feel at the end of the dialogue.} \\
\\ 
\texttt{Trial \#2: New plan: } \\
\\ 
\texttt{Upon re{-}examining the dialogue, I realized that I focused too much on the surface{-}level emotions expressed by Felix and Vanessa, but didn't consider the underlying dynamics of their relationship and the severity of Felix's actions. } \\
\\ 
\texttt{To improve my prediction, I will pay closer attention to the following:} \\
\\ 
\texttt{1. The extent of Felix's betrayal: I will consider the impact of Felix's actions on Vanessa's reputation, career, and trust in their friendship.} \\
\texttt{2. Vanessa's emotional state: I will analyze her reactions and language to gauge the depth of her hurt, anger, and disappointment.} \\
\texttt{3. Felix's lack of accountability: I will assess his responses to see if he truly understands the gravity of his mistakes and takes responsibility for his actions.} \\
\\ 
\texttt{With these factors in mind, I will reassess the emotions Felix is likely to feel at the end of the dialogue.} \\
\\ 
\texttt{Relief: 0} \\
\texttt{Despair: 9} \\
\texttt{Vindication: 0} \\
\texttt{Disappointment: 9} \\
\\ 
\texttt{Trial \#3: New plan: } \\
\\ 
\texttt{Upon re{-}examining the dialogue, I realized that I focused too much on the surface{-}level emotions expressed by Felix and Vanessa, but didn't consider the underlying dynamics of their relationship and the severity of Felix's actions. } \\
\\ 
\texttt{To improve my prediction, I will pay closer attention to the following:} \\
\\ 
\texttt{1. The extent of Felix's betrayal: I will consider the impact of Felix's actions on Vanessa's reputation, career, and trust in their friendship.} \\
\texttt{2. Vanessa's emotional state: I will analyze her reactions and language to gauge the depth of her hurt, anger, and disappointment.} \\
\texttt{3. Felix's lack of accountability: I will assess his responses to see if he truly understands the gravity of his mistakes and takes responsibility for his actions.} \\
\\ 
\texttt{However, I also noticed that I gave high scores to both Despair and Disappointment, which might be overlapping emotions. To refine my prediction, I will consider the possibility that Felix is experiencing a more intense sense of regret and remorse, which could be a combination of these two emotions.} \\
\\ 
\texttt{With these factors in mind, I will reassess the emotions Felix is likely to feel at the end of the dialogue.} \\
\\ 
\texttt{Relief: 0} \\
\texttt{Despair: 8} \\
\texttt{Vindication: 0} \\
\texttt{Disappointment: 8} \\
\\ 
\\ 
\\ 
\texttt{Your task is to predict the likely emotional responses of a character in this dialogue:} \\
\\ 
\texttt{Felix: Vanessa, I...} \\
\texttt{Vanessa: Save it, Felix. I don't want to hear your excuses.} \\
\texttt{Felix: It's not an excuse, it's an apology.} \\
\texttt{Vanessa: And you think that makes it better? } \\
\texttt{Felix: I'm not trying to make it better. I just... I messed up, alright?} \\
\texttt{Vanessa: You "messed up"? You spread lies about me, Felix. You ruined my reputation, my career!} \\
\texttt{Felix: I know, and I'm sorry. It was a stupid thing to do. I was... I was just...} \\
\texttt{Vanessa: What? Bored? Looking for a thrill? } \\
\texttt{Felix: I... I don't know.} \\
\texttt{Vanessa: You don't know. Of course, you don't.} \\
\texttt{Felix: Vanessa, I...} \\
\texttt{Vanessa: Stop. Just stop, Felix. I don't want your apologies. I don't want your excuses. I don't want anything from you.} \\
\texttt{Felix: I... I understand.} \\
\texttt{Vanessa: Do you? Because I don't think you do. I don't think you understand how much you hurt me. How much you betrayed me.} \\
\texttt{Felix: I... I didn't mean to...} \\
\texttt{Vanessa: But you did, Felix. You did. And now, I can't trust you anymore. I can't be your friend.} \\
\texttt{{[}End dialogue{]}} \\
\\ 
\texttt{At the end of this dialogue, Felix would feel...} \\
\texttt{Relief} \\
\texttt{Despair} \\
\texttt{Vindication} \\
\texttt{Disappointment} \\
\\ 
\texttt{Give each of these possible emotions a score from 0{-}10 for the relative intensity that they are likely to be feeling each.} \\
\\ 
\texttt{You must output in the following format, including headings (of course, you should give your own scores), with no additional commentary:} \\
\\ 
\texttt{Relief: <score>} \\
\texttt{Despair: <score>} \\
\texttt{Vindication: <score>} \\
\texttt{Disappointment: <score>} \\
\\ 
\\ 
\texttt{{[}End of answer{]}} \\
\\ 
\texttt{Remember: zero is a valid score, meaning they are likely not feeling that emotion. You must score at least one emotion > 0.} \\
\\ 
\texttt{Your answer:} \\
            }
        }
    }
     \\ \\

    \theutterance \stepcounter{utterance}  
    & \multicolumn{4}{p{0.6\linewidth}}{
        \cellcolor[rgb]{0.561,0.737,0.561}{
            \makecell[{{p{\linewidth}}}]{
                \texttt{\tiny{[P1$\rangle$GM]}}
                \texttt{Relief: 0} \\
\texttt{Despair: 9} \\
\texttt{Vindication: 0} \\
\texttt{Disappointment: 10} \\
            }
        }
    }
    & & \\ \\

    \theutterance \stepcounter{utterance}  
    & & & \multicolumn{2}{p{0.3\linewidth}}{
        \cellcolor[rgb]{0.9,0.9,0.9}{
            \makecell[{{p{\linewidth}}}]{
                \texttt{\tiny{[GM$|$GM]}}
                \texttt{Relief: 0} \\
\texttt{Despair: 9} \\
\texttt{Vindication: 0} \\
\texttt{Disappointment: 10} \\
            }
        }
    }
    & & \\ \\

    \theutterance \stepcounter{utterance}  
    & & & \multicolumn{2}{p{0.3\linewidth}}{
        \cellcolor[rgb]{0.9,0.9,0.9}{
            \makecell[{{p{\linewidth}}}]{
                \texttt{\tiny{[GM$|$GM]}}
                \texttt{Relief: 1} \\
\texttt{Despair: 8} \\
\texttt{Vindication: 0} \\
\texttt{Disappointment: 9} \\
            }
        }
    }
    & & \\ \\

    \theutterance \stepcounter{utterance}  
    & & & \multicolumn{2}{p{0.3\linewidth}}{
        \cellcolor[rgb]{0.9,0.9,0.9}{
            \makecell[{{p{\linewidth}}}]{
                \texttt{\tiny{[GM$|$GM]}}
                \texttt{game\_result = LOSE} \\
            }
        }
    }
    & & \\ \\

    \theutterance \stepcounter{utterance}  
    & & \multicolumn{4}{p{0.6\linewidth}}{
        \cellcolor[rgb]{0.95,0.95,0.95}{
            \makecell[{{p{\linewidth}}}]{
                \texttt{\tiny{[P1$\leftrightarrow$P]}}
                \texttt{Post{-}episode Reflexion (not sent to GM):} \\
\texttt{New plan: } \\
\\ 
\texttt{Upon re{-}examining the dialogue, I realized that I focused too much on the surface{-}level emotions expressed by Felix and Vanessa, but didn't consider the underlying dynamics of their relationship and the severity of Felix's actions. } \\
\\ 
\texttt{To improve my prediction, I will pay closer attention to the following:} \\
\\ 
\texttt{1. The extent of Felix's betrayal: I will consider the impact of Felix's actions on Vanessa's reputation, career, and trust in their friendship.} \\
\texttt{2. Vanessa's emotional state: I will analyze her reactions and language to gauge the depth of her hurt, anger, and disappointment.} \\
\texttt{3. Felix's lack of accountability: I will assess his responses to see if he truly understands the gravity of his mistakes and takes responsibility for his actions.} \\
\\ 
\texttt{However, I also noticed that I gave high scores to both Despair and Disappointment, which might be overlapping emotions. To refine my prediction, I will consider the possibility that Felix is experiencing a more intense sense of regret and remorse, which could be a combination of these two emotions.} \\
\\ 
\texttt{Additionally, I will also consider the fact that Felix's attempts to apologize and make excuses are met with rejection and anger from Vanessa. This could indicate that Felix is not only feeling regret and remorse but also a sense of hopelessness and despair about the state of their relationship.} \\
\\ 
\texttt{With these factors in mind, I will reassess the emotions Felix is likely to feel at the end of the dialogue.} \\
\\ 
\texttt{Relief: 0} \\
\texttt{Despair: 9} \\
\texttt{Vindication: 0} \\
\texttt{Disappointment: 9} \\
            }
        }
    }
    & \\ \\

\end{supertabular}
}

\end{document}
